\chapter{Revisão de Literatura  }
\label{cap:revisao-literatura}

% \section{Aplicações em Redes Veiculares}
% \label{sec:aplicacoes-em-redes-veiculares}

% \section{Técnicas de Offloading em Ambientes Veiculares}
% \label{sec:tecnicas-de-offloading-em-ambientes-veiculares}

% \section{Modelos Híbridos Fuzzy e Aprendizado por Reforço}
% \label{sec:modelos-hibridos-fuzzy-e-aprendizado-por-reforco}

% \section{Problemas e Desafios Atuais}
% \label{sec:problemas-e-desafios-atuais}

Este capítulo apresenta uma análise crítica e comparativa dos trabalhos mais recentes e relevantes para o desenvolvimento desta proposta. A literatura é organizada segundo uma taxonomia adaptada de~\cite{de2020computation}, estruturada em duas vertentes. A Seção~\ref{sec:padrao-comunicacao} examina o Padrão de Comunicação adotado pelos trabalhos relacionados, considerando as tecnologias de rede empregadas e os servidores de borda utilizados; e a Seção~\ref{sec:problema} analisa a Formulação do Problema de descarregamento, com ênfase nas estratégias baseadas em \gls{DRL}, na engenharia de estado e recompensa, e nos ambientes de simulação adotados. Por fim, a Seção~\ref{sec:consideracoes-finais-revisao} sintetiza as lacunas identificadas e posiciona a contribuição desta proposta frente ao estado da arte.

\section*{Taxonomia}

Para classificar e relacionar o estado da arte das soluções de descarregamento de tarefas em ambientes de \gls{VEC}, adotou-se uma versão adaptada da taxonomia proposta em~\cite{de2020computation}, sendo ela dividida em duas vertentes estruturais: Padrão de Comunicação, que avalia as tecnologias de rede e os servidores utilizados; e Problema, que analisa as estratégias com ênfase em modelos \gls{DRL}. Dessa maneira, os artigos relacionados foram classificados conforme a Figura~\ref{fig:taxonomia_adaptada}.

\begin{figure}[ht!]
\centering
\resizebox{\linewidth}{!}{%
\begin{forest}
  for tree={
    font=\sffamily\normalsize,
    draw=teal!80!black,
    very thick,
    rounded corners=4pt,
    fill=teal!70!black,
    text=white,
    align=center,
    edge={draw=teal!60!black, thick},
    l sep=6mm,
    s sep=2mm,
    if level=0{
      minimum height=8mm,
      font=\sffamily\normalsize\bfseries,
    }{
      if level=1{
        fill=teal!60!black,
        font=\sffamily\normalsize\bfseries,
        minimum height=7mm,
      }{
        if level=2{
          fill=teal!50!black,
          font=\sffamily\normalsize\bfseries,
          minimum height=6mm,
        }{
          % Leaf nodes
          fill=teal!40!black,
          font=\sffamily\normalsize,
          draw=teal!50!black,
          rounded corners=3pt,
          align=left,
          inner sep=2pt,
          thin,
        }
      }
    }
  }
  [Descarregamento\\Computacional Veicular
    [Padrão de\\Comunicação
      [Tecnologia
        [{$\triangleright$ C-\gls{V2I} / C-\gls{V2X}\\$\triangleright$ 5G / 6G}]
      ]
      [Servidor
        [{$\triangleright$ \gls{MEC} (Borda Fixa)\\$\triangleright$ Nuvens Veiculares\\$\triangleright$ Nuvem}]
      ]
    ]
    [Problema
      [Estratégia
        [Função Objetivo
          [{$-$ Tempo de\\~~~Resposta\\$-$ Consumo\\~~~Energético\\$-$ Custos Financeiros\\$+$ Estabilidade de\\~~~Filas (Lyapunov)\\$+$ Taxa de Sucesso\\$+$ Nível de Confiança}]
        ]
        [Eng. de Estado
          [{$\triangleright$ Tamanho das Filas\\$\triangleright$ Prioridade / \textit{Deadline}\\$\triangleright$ Escores de Confiança\\$\triangleright$ Ângulos e Orientação}]
        ]
        [Eng. de Recompensa
          [{$\triangleright$ Pesos Dinâmicos\\$\triangleright$ Penalização Severa\\$\triangleright$ Recompensa Global\\$\triangleright$ Estágios Sequenciais\\$\triangleright$ Drift-plus-penalty}]
        ]
        [Modelos
          [{$\triangleright$ \gls{DQN}, D3QN, P-\gls{DQN}\\$\triangleright$ \gls{SAC}, DDPG, QDPG\\$\triangleright$ L-MADDPG, SIDDPG\\$\triangleright$ Fusões: FL, DT, GNN}]
        ]
      ]
      [Ambiente\\Simulado
        [{$\triangleright$ Mobilidade: SUMO,\\~~~PTV Vissim\\$\triangleright$ Rede: ns-3\\$\triangleright$ Agente: OpenAI Gym\\$\triangleright$ Dados Reais (Traces)}]
      ]
    ]
  ]
\end{forest}
}
\Caption{\label{fig:taxonomia_adaptada} Taxonomia adaptada para classificação dos trabalhos relacionados em descarregamento computacional veicular, com ênfase na modelagem de \gls{DRL} e ambiente de simulação.}
\Fonte{Adaptado de~\citeonline{de2020computation}.}
\end{figure}

\section{Padrão de Comunicação}
\label{sec:padrao-comunicacao}
No descarregamento de tarefas, o cenário de comunicação celular vem ganhando bastante destaque. Neste quesito, suas tecnologias vêm evoluindo rapidamente e ditam a maioria das pesquisas aqui apresentadas. A seguir são apresentados alguns trabalhos que utilizam tecnologias \gls{5G} e além (\gls{B5G}/\gls{6G}) para atender às exigentes demandas relacionadas a aplicações veiculares.

\subsection{Tecnologia}
A evolução das tecnologias de comunicação veicular destaca a transição de padrões tradicionais, como o \gls{DSRC} e o IEEE 802.11bd, para a comunicação celular aliada à \gls{IA}. Essa mudança visa, sobretudo, possibilitar baixas latências e alta vazão de dados em cenários de alta mobilidade, contribuindo para resolver o problema de descarregamento de tarefas. Nesse contexto, diversos trabalhos, como os de~\cite{Moghaddasi2023, sidwini, li2025, meghamala2025}, utilizam interfaces \gls{C-V2X} e arquiteturas de borda para o compartilhamento de informações em aplicações veiculares avançadas. 

O estudo de~\cite{Moghaddasi2023} propõe um algoritmo de \gls{DRL}, especificamente o \gls{DDQN}, em redes \gls{5G} para solucionar o problema de descarregamento de tarefas utilizando a comunicação \gls{V2I} para alcançar eficiência energética e a diminuição da latência. O trabalho de~\cite{sidwini} explora o uso de algoritmos avançados de \gls{DRL}, como o \gls{SAC}, para otimizar o descarregamento dinâmico entre dispositivos locais, servidores de borda e a nuvem, equilibrando o consumo de energia e o atraso na execução. Em um cenário direcionado aos sistemas \gls{6G},~\cite{li2025} introduz um esquema inteligente baseado em \gls{DQN} em redes \gls{V2X} assistidas por \gls{MEC} para otimizar conjuntamente a associação aos servidores, a taxa de descarregamento de tarefas e a aceleração dos veículos. Essa abordagem visa minimizar o atraso do sistema e evitar colisões de trânsito, utilizando uma função de recompensa segmentada que respeita restrições de atraso, energia e segurança. Por fim,~\cite{meghamala2025} propõe um \textit{framework} de \gls{MARL} para redes \gls{5G} que combina o algoritmo \gls{SAC} com arquiteturas hierárquicas (\gls{MATD3} e \gls{TD2PG}). Essa solução visa a alocação dinâmica de subcanais, potência de transmissão e recursos computacionais em servidores \gls{MEC}, levando em consideração a mobilidade dos usuários e o isolamento das fatias de rede (\textit{network slicing}).

\subsection{Servidor}
O \gls{VEC} é um paradigma amplo que integra dois ambientes complementares de execução remota: a infraestrutura de borda fixa, composta por servidores \gls{MEC} acoplados a \glspl{RSU} ou estações base e acessíveis via comunicação \gls{V2I}; e as Nuvens Veiculares, formadas pelos recursos computacionais ociosos dos próprios veículos e acessíveis por meio de canais diretos \gls{V2V}. Esses dois ambientes podem operar de forma isolada ou integrada. Adicionalmente, quando as restrições de latência permitem, tarefas podem ser delegadas a uma infraestrutura de nuvem remota, que oferece maior capacidade de processamento ao custo de latência consideravelmente superior. A seguir, são descritas as abordagens modeladas para cada um desses ambientes de execução nos trabalhos relacionados.

\section*{\texorpdfstring{\gls{MEC}}{\gls{MEC}}}
Servidores \gls{MEC} contribuem significativamente para atender exigências de \gls{URLLC} em \gls{5G} \gls{NR}. Eles ficam localizados na borda da rede e podem ser constituídos de equipamentos robustos, capazes de auxiliar no processamento pesado de tarefas delegadas por outros dispositivos com recursos limitados.

Nesse contexto, esses servidores frequentemente estão conectados a \gls{RSU} e estações base para oferecer suporte computacional de baixa latência aos veículos. Aproveitando essa infraestrutura de borda, o estudo de~\cite{zhao2025} propõe o esquema \gls{PDQKM} para otimizar o descarregamento de múltiplas tarefas, permitindo reduzir a latência total do sistema. Por outro lado, o trabalho de~\cite{wu2025} utiliza \gls{GAN} para expandir amostras de tarefas evitando o desequilíbrio de dados, criando um modelo de \gls{FL} para realizar o descarregamento de tarefas de forma energeticamente eficiente em servidores \gls{MEC}. O estudo de~\cite{mohammad2025} apresenta o \textit{framework} EMEC-IoVTOF, que integra \gls{DRL} com a técnica de \gls{PSO} para otimizar a alocação de recursos em sistemas \gls{IoV}. A solução calcula custos de descarregamento para evitar a convergência para ótimos locais, conseguindo reduções significativas de latência e consumo de energia, além de contornar restrições de largura de banda em cenários veiculares dinâmicos. Por sua vez, o estudo de~\cite{liu2025} utiliza o método de aproximação de Bernstein para converter restrições probabilísticas de interferência em formato tratável, minimizando uma função de utilidade que equilibra atrasos na transmissão e na computação para o \gls{C-MEC}.

Nessa mesma direção, o trabalho de~\cite{vieira2025} propõe a política de descarregamento de tarefas \gls{FOFF}, voltada para redes \gls{VEC} habilitadas por \gls{5G} via comunicação \gls{V2I}. Em vez de recorrer a técnicas iterativas de aprendizado, o FOFF emprega o método de \gls{MCDM} \textit{Fuzzy}-\gls{TOPSIS} com Números Fuzzy Triangulares (TFNs) para capturar a incerteza e a variabilidade dos recursos dos servidores \gls{MEC} em tempo de decisão, sem a necessidade de uma fase de treinamento prévio. A política resultante determina, para cada tarefa gerada pelo veículo, se o processamento deve ocorrer localmente na \gls{OBU} ou ser descarregado para o servidor de borda fixo, com base na pontuação de proximidade relativa calculada pelo \textit{Fuzzy}-\gls{TOPSIS} em função do tempo de resposta e do consumo energético. Avaliado em um simulador com heterogeneidade computacional, o \gls{FOFF} demonstrou superioridade sobre métodos gulosos (GTT), heurísticas particionadas (FP-Sched) e abordagens baseadas em algoritmos genéticos (HOFF), alcançando reduções no consumo energético de até 72,11\%.

\section*{Nuvens Veiculares}

As Nuvens Veiculares constituem o ambiente de execução cooperativo do paradigma \gls{VEC} baseado em comunicação \gls{V2V}. Nesse ambiente, os recursos computacionais ociosos das unidades embarcadas (\gls{OBU}) dos próprios veículos são disponibilizados para auxiliar nós vizinhos por meio de canais de comunicação direta, sem necessidade de intermediação por infraestruturas estáticas da pista. O estudo de~\cite{xie2025} propõe utilizar \gls{RIS} juntamente com \gls{VEC} para reduzir custos de implantação, ao mesmo tempo que busca maximizar \gls{SOE} via algoritmo \gls{PPO}, otimizando conjuntamente decisões de descarregamento, alocação de recursos computacionais e vetores de mudança de fase. Por sua vez,~\cite{wu2025ava} utiliza \gls{FL} e \gls{MADDPG}, empregando mecanismos de incentivo baseados em reputação para determinar a proporção ideal de recursos cedidos por veículos vizinhos, visando minimizar latência e consumo de energia e garantindo alta taxa de conclusão em tarefas computacionalmente intensivas. O trabalho de~\cite{long2025} desenvolve o algoritmo \gls{CTSRAO}, que integra dispositivos finais, borda veicular e nuvem para realizar descarregamento parcial de tarefas, separando o problema em subproblemas convexos para otimizar de forma colaborativa o \textit{task splitting} (taxa de divisão de tarefas) e a alocação de recursos computacionais, alcançando equilíbrio de desempenho global do sistema. Em~\cite{abbas2025}, propõe-se um algoritmo baseado em aprendizado descentralizado colaborativo, que orquestra a execução concorrente de tarefas fragmentadas pelos veículos presentes na área de cobertura, otimizando conjuntamente a alocação de canais sem fio com os recursos computacionais. Dessa forma, minimiza-se a latência total e o consumo de energia em ambientes veiculares dinâmicos.

\section*{Discussão}
\begin{table}[h]
    \centering
    \Caption{\label{tab:padrao_comunicacao}Resumo dos aspectos de Padrão de Comunicação, Servidor e Condições de Ambiente dos trabalhos relacionados}
    \renewcommand{\arraystretch}{1.2}
    \small
    \setlength{\tabcolsep}{2pt}
    \begin{tabularx}{\textwidth}{|>{\raggedright\arraybackslash}X|c|c|c|c|c|c|c|c|}
        \hline
        \multirow{2}{*}{Referência} & \multicolumn{2}{c|}{Tecnologia} & \multicolumn{3}{c|}{Servidor} & \shortstack{Observabilidade\\Parcial} & \multirow{2}{*}{\shortstack{Canal\\Perfeito}} & \shortstack{Carga de\\Trabalho} \\ \cline{2-6}
                                             & \gls{V2I} & \shortstack{\textit{Sidelink}\\(\gls{V2V})} & \gls{MEC} & \shortstack{Nuvens\\Veiculares} & Nuvem & \shortstack{(Incerteza\\de Estado)} & & \shortstack{(Diferentes\\regimes)} \\ \hline
       ~\cite{Moghaddasi2023} & \checkmark & & \checkmark & & & & \checkmark & \checkmark \\ \hline
       ~\cite{sidwini}           & \checkmark & & \checkmark & & \checkmark & & \checkmark & \\ \hline
       ~\cite{li2025}                 & & \checkmark & \checkmark & & & & \checkmark & \\ \hline
       ~\cite{meghamala2025}   & \checkmark & & \checkmark & & & & \checkmark & \\ \hline
       ~\cite{zhao2025}             & \checkmark & & \checkmark & & & & \checkmark & \\ \hline
       ~\cite{wu2025}                 & \checkmark & & \checkmark & & & & \checkmark & \\ \hline
       ~\cite{mohammad2025}     & & \checkmark & \checkmark & & & & \checkmark & \\ \hline
       ~\cite{liu2025}               & \checkmark & & \checkmark & & \checkmark & & & \\ \hline
       ~\cite{xie2025}               & \checkmark & & & \checkmark & & & \checkmark & \\ \hline
       ~\cite{wu2025ava}              & & \checkmark & & \checkmark & & & \checkmark & \\ \hline
       ~\cite{long2025}             & \checkmark & \checkmark & & \checkmark & \checkmark & & \checkmark & \\ \hline
       ~\cite{abbas2025}           & \checkmark & & & \checkmark & & & \checkmark & \\ \hline
       ~\cite{vieira2025}      & \checkmark & & \checkmark & & & \checkmark & \checkmark & \\ \hline
    \end{tabularx}
    \Fonte{Elaborado pelo autor}
\end{table}

Analisando a Tabela~\ref{tab:padrao_comunicacao}, constata-se forte tendência de padronização de tecnologias celulares avançadas para comunicação veicular. Grande parte dos trabalhos recentes analisados concentra-se na utilização de conectividade \gls{V2I} via redes \gls{5G} e \gls{6G}. Em relação à comunicação entre os nós, embora o padrão \gls{C-V2X} englobe diferentes interfaces, diversos trabalhos referem-se estritamente à utilização das interfaces \textit{Sidelink} destas comunicações para troca direta de dados (\gls{V2V}). Já em relação à infraestrutura, a adoção de servidores \gls{MEC} centralizados em unidades \gls{RSU} ainda é predominante, mas observa-se tendência crescente de pesquisas que exploram o ambiente de Nuvens Veiculares, constituinte do paradigma \gls{VEC}, tratando os próprios veículos como provedores de recursos. Trabalhos como os de~\cite{liu2025} e~\cite{long2025} adotam arquiteturas colaborativas integrando \gls{MEC} com nuvem e Nuvens Veiculares com nuvem, respectivamente.

Outro aspecto identificado refere-se ao tratamento da observabilidade parcial, caracterizada na tabela pela coluna "Incerteza de Estado". Com exceção de~\citeonline{vieira2025}, em que Números Fuzzy Triangulares (TFNs) são empregados para modelar explicitamente a imprecisão dos recursos disponíveis nos servidores \gls{MEC}, os demais trabalhos identificados modelam o ambiente assumindo conhecimento perfeito do cenário, como a disponibilidade de \gls{CPU} dos servidores no exato momento da tomada de decisão. Dessa forma, agentes independentes podem decidir descarregar tarefas para um mesmo servidor simultaneamente com base em informações desatualizadas. Além disso, não há menção nos demais trabalhos sobre imperfeições do canal de comunicação. Portanto, assumir tais características cria políticas de decisão pouco robustas frente à dinâmica real.

Outra limitação identificada diz respeito à modelagem de cargas de trabalho. Com exceção do estudo de~\cite{Moghaddasi2023}, os demais trabalhos avaliam seus algoritmos sob regimes de tráfego constantes ou com taxa de chegada fixa por episódio. Em cenários de tráfego reais, eventos imprevisíveis podem causar rajadas instantâneas (\textit{bursts}) no envio de pacotes críticos, alterando a distribuição de chegada das tarefas. Modelos \gls{DRL} treinados exclusivamente sem variação dinâmica de carga tendem a sofrer degradação de desempenho com tais mudanças de regime.

Nesse contexto, o modelo proposto neste trabalho visa endereçar essas lacunas ao modelar, desde a concepção do ambiente, tais dificuldades e estocasticidades inerentes a redes veiculares reais, conforme a adoção das tecnologias supramencionadas.

\section{Problema}
\label{sec:problema}
Com a infraestrutura de comunicação estabilizada, a escolha de abordagens apropriadas constituem um passo crítico para a garantia dos requisitos no ambiente veicular. A formulação de soluções para o descarregamento de tarefas dinâmico e em tempo real apresenta elevada complexidade. Para contornar essa dificuldade, diversos trabalhos na literatura optam por simplificar o escopo, desenvolvendo abordagens focadas em uma única métrica, como a minimização exclusiva da latência ou do consumo de energia~\cite{feng2017ave, hou2016vehicular, ren2018latency, ren2019collaborative, zhang2019delay}. No entanto, o estudo de~\cite{qin2020} enfatiza que otimizar um único objetivo isoladamente não reflete a realidade das redes, enquanto~\cite{zhang2024} reitera que a busca por um compromisso dinâmico entre tempo de resposta e consumo energético transforma a alocação em um problema incerto de otimização multiobjetivo NP-difícil. Diante dessa constatação, observa-se a necessidade do desenvolvimento de políticas que considerem múltiplos critérios simultaneamente, com o intuito de projetar propostas robustas e capazes de atender às demandas operacionais e funcionais de longo prazo.

A modelagem de um cenário realista precisa lidar fundamentalmente com incertezas espaciais da rede. A alta velocidade dos veículos altera significativamente a infraestrutura veicular, mascarando os estados reais e comprometendo a acurácia dos dados divulgados pelos nós de computação~\cite{liu2025}. A somatória entre a imprevisibilidade de processos no interior da borda~\cite{abbas2025} e o natural decréscimo de precisão proveniente do atraso por telemetria exige uma abordagem maleável. Decisões sustentadas integralmente por lógicas determinísticas tornam-se insatisfatórias, contribuindo para que os modelos não atendam adequadamente os requisitos de qualidade no descarregamento de tarefas~\cite{kandali2025, zhao2025}.

\subsection{Estratégia}
Para contornar grande parte dos vieses supracitados, a literatura contemporânea aponta uma considerável substituição do uso de heurísticas e meta-heurísticas por soluções mais sofisticadas e adaptativas. A respeito desta transformação, destaca-se a utilização de \gls{DRL}, por conta de sua capacidade de aprimorar sequências de ações de longo prazo extraídas através de intensas interações com ambiente simulado.

Antes da consolidação do \gls{DRL} como paradigma predominante, abordagens de \gls{MCDM} foram investigadas como alternativa para lidar com a natureza multiobjetivo e incerta do problema de descarregamento. O trabalho de~\cite{vieira2025} demonstrou a viabilidade do método \textit{Fuzzy}-\gls{TOPSIS} operando com Números Fuzzy Triangulares (TFNs) como mecanismo de decisão eficiente: sem incorrer em fase de treinamento, a solução computa pontuações de proximidade relativa para cada alternativa de execução --- local (\gls{OBU}) ou remota (\gls{MEC}) --- e seleciona a que melhor atende ao compromisso entre latência e consumo energético sob condições de incerteza. Os pesos dos critérios de decisão pertencem ao intervalo \(\left[0, 1\right]\) e refletem a importância relativa de cada métrica de desempenho, enquanto os TFNs capturam a imprecisão inerente às medições dos recursos dos servidores de borda. Contudo, essa classe de soluções apresenta uma limitação estrutural relevante: tanto os pesos dos critérios quanto os conjuntos fuzzy associados são pré-configurados antes da implantação e permanecem inalterados durante toda a execução. Quando o regime de tráfego sofre uma mudança abrupta --- por exemplo, devida a rajadas (\textit{bursts}) de requisições críticas ou à degradação súbita da disponibilidade de servidores --- a política \textit{Fuzzy}-\gls{TOPSIS} não dispõe de mecanismo autônomo para recalibrar seus parâmetros, o que compromete a qualidade das decisões nesses cenários imprevistos. Essa rigidez constitui a motivação central da evolução metodológica proposta neste trabalho: ao integrar o raciocínio \textit{fuzzy}, responsável pela quantificação da incerteza pontual, com a capacidade de exploração contínua e atualização de política proporcionada pelo algoritmo \gls{SAC} de \gls{DRL}, obtém-se um agente capaz de se adaptar a novos regimes de tráfego sem re-parametrização manual, combinando a eficiência multicritério de um com a plasticidade adaptativa do outro.

De fato, o escopo científico atual tem explorado ativamente o uso de diversas estratégias  integradas ao \gls{DRL}. Uma delas é a integração com aprendizado \gls{FL}, que permite o treinamento colaborativo de modelos \gls{DRL} entre múltiplos agentes, preservando a privacidade dos dados locais~\cite{chen2025}. Outra frente mescla esquemas \gls{DRL} com tecnologia \gls{DT}, possibilitando que soluções sejam testadas e simuladas virtualmente antes da execução no ambiente físico~\cite{wei2024}. Para lidar com o vasto espaço de estados e com cenários altamente variáveis, pesquisas recentes fundem arquiteturas \gls{DRL} com redes \gls{GNN} para capturar de forma eficiente as relações e as dependências espaciais entre os nós veiculares~\cite{wang2024, tang2024}. Adicionalmente, métodos \gls{TL} têm sido utilizados para acelerar a adaptação de veículos recém-conectados em infraestruturas baseadas em \gls{SDN}, em que a separação entre os planos de controle e de dados facilita a gerência global~\cite{baker2024}. Por fim, destaca-se a incorporação da teoria de otimização de Lyapunov acoplada a métodos \gls{DRL} para garantir matematicamente a estabilidade contínua das filas de tarefas a longo prazo~\cite{kumar2023}.

\section*{Função Objetivo}
Além do compromisso entre a minimização da latência de processamento e a redução do consumo de energia, a literatura recente incorpora a otimização de custos operacionais e financeiros da rede~\cite{ma2025}. A estabilidade das filas de tarefas (\textit{queue stability}) a longo prazo também é um objetivo considerado, frequentemente abordado por meio da otimização de Lyapunov acoplada a técnicas de aprendizado~\cite{kumar2023}. Outros estudos buscam maximizar a taxa de sucesso das tarefas sob restrições de prioridade e da capacidade dos nós de computação~\cite{kumari2025}. Adicionalmente, o gerenciamento de confiança (\textit{trust management}) dos nós de borda desponta como um objetivo crítico em infraestruturas \gls{SDN} para evitar descarregamentos falhos ou maliciosos~\cite{baker2024}.

\section*{Engenharia de Estado e Recompensa}
No \gls{MDP}, a engenharia do espaço de estado não busca prever matematicamente as transições do sistema, uma vez que algoritmos \gls{DRL} se destacam justamente por gerenciar o controle em ambientes estocásticos e dinâmicos sem conhecimento prévio dessas transições~\cite{fan2024}. Em vez disso, o estado é projetado para permitir que o agente avalie diretamente a dinâmica da rede com base nos resultados observados. Para isso, o espaço de estado deve utilizar métricas além das básicas, incorporando variáveis que capturem a dinamicidade dos nós de computação~\cite{kumar2023} ou até mesmo escores de confiança dos dispositivos candidatos~\cite{baker2024}. A literatura recente expande ainda mais essa abstração ao integrar as características das tarefas --- englobando requisitos de recursos, níveis de prioridade e prazos máximos toleráveis (\textit{deadlines})~\cite{fan2024, wei2024}. Além disso, para refinar a percepção física, as pesquisas incluem o ângulo de desvio e a orientação direcional dos veículos em relação aos servidores. Essa consciência espacial de granularidade fina complementa o estado atual, otimizando a escolha do nó alvo de forma proativa~\cite{farimani2024}.

A respeito da função de recompensa, é uma boa prática~\cite{sutton2018} realizar o dimensionamento adequado (\textit{scaling}) das métricas e a normalização dos dados para evitar que uma única variável domine o processo de convergência. Além disso, alguns trabalhos utilizam pesos dinâmicos em vez de estáticos, ajustando-os em tempo real de acordo com as condições momentâneas da rede e com o nível de urgência das tarefas~\cite{kumari2025, fan2024, shi2023}. Já a penalização para quando um agente escolhe uma ação ruim --- como a violação do tempo máximo tolerável (\textit{deadline}) ou quebras de comunicação por evasão de cobertura --- é integrada à recompensa para guiar o agente a priorizar o cumprimento dos requisitos de latência e confiabilidade. Também se observa uma forte transição para mecanismos de recompensa globais, nos quais os agentes colaboram buscando a eficiência sistêmica da rede, mesmo que em detrimento do custo ótimo individual imediato~\cite{farimani2024}. Em cenários de descarregamento parcial, a recompensa é formulada dividindo a ação em múltiplos estágios sequenciais, avaliando o impacto do particionamento da tarefa de forma independente da alocação de recursos~\cite{tian2025}. Por fim, a teoria de Lyapunov é frequentemente utilizada na recompensa por meio de funções \textit{drift-plus-penalty}, garantindo matematicamente que o agente seja penalizado caso a redução excessiva do consumo energético coloque em risco a estabilidade das filas do sistema a longo prazo~\cite{kumar2023}.

\section*{Modelos}
No que tange aos modelos, existem diversos trabalhos focados em otimizar a forma como o agente observa o estado e seleciona ações eficientemente. Variantes baseadas em valor, como redes \gls{DQN} aprimoradas, realizam a seleção de servidores reconhecendo que nem todas as tarefas computacionais geradas possuem a mesma importância ou urgência~\cite{kumari2025}. Para espaços de ação contínuos, arquiteturas Ator-Crítico dominam o cenário. O algoritmo \gls{SAC} é empregado para evitar a convergência prematura para ótimos locais, utilizando a maximização da entropia estocástica para explorar melhor as opções de descarregamento e \textit{caching}~\cite{bao2025}. Para lidar com o problema de espaços de ação híbridos --- nos quais decisões discretas de seleção de servidor são combinadas com a alocação contínua de recursos computacionais ---, alguns trabalhos propõem métodos como o algoritmo \gls{P-DQN}~\cite{ma2025} ou mesmo fusões entre modelos, como o \gls{QDPG}, que une as arquiteturas \gls{DDPG} e \gls{D3QN} aliadas a algoritmos de clusterização \textit{K-Means}~\cite{wei2024}. Em sistemas compostos por múltiplos agentes, o modelo \gls{L-MADDPG} utiliza a otimização de Lyapunov para garantir matematicamente a estabilidade do sistema a longo prazo~\cite{kumar2023}, enquanto abordagens como o modelo \gls{SIDDPG} iteram múltiplos estados internamente para refinar o particionamento de sub-tarefas de forma sequencial~\cite{tian2025}. A Tabela~\ref{tab:drl_action_spaces} mapeia alguns algoritmos \gls{DRL} ao tipo de espaço de ação adequado.

\begin{table}[ht!]
    \centering
    \Caption{\label{tab:drl_action_spaces}Classificação dos algoritmos clássicos de \gls{DRL} quanto à adequação ao espaço de ação}
    \renewcommand{\arraystretch}{1.3}
    \setlength{\tabcolsep}{2pt}
    \begin{tabularx}{\textwidth}{|>{\raggedright\arraybackslash}X|c|>{\raggedright\arraybackslash}X|}
        \hline
        Algoritmo \gls{DRL} & Espaço de Ação & Referência Base \\ \hline
        \gls{DQN} (\textit{Deep Q-Network}) & Discreto &~\cite{uddin2025, farimani2024} \\ \hline
        Double \gls{DQN} & Discreto &~\cite{uddin2025} \\ \hline
        Dueling \gls{DQN} & Discreto &~\cite{uddin2025} \\ \hline
        \gls{DDPG} (\textit{Deep Deterministic Policy Gradient}) & Contínuo &~\cite{uddin2025, farimani2024} \\ \hline
        \gls{TD3} (\textit{Twin Delayed DDPG}) & Contínuo &~\cite{uddin2025} \\ \hline
        \gls{PPO} (\textit{Proximal Policy Optimization}) & Discreto e Contínuo &~\cite{uddin2025} \\ \hline
        TRPO (\textit{Trust Region Policy Optimization}) & Discreto e Contínuo &~\cite{uddin2025} \\ \hline
        \gls{SAC} (Soft \textit{Actor}-\textit{Critic}) & Contínuo &~\cite{uddin2025, farimani2024} \\ \hline
    \end{tabularx}
    \Fonte{Elaborado pelo autor}
\end{table}

\section*{Ambiente Simulado}
O treinamento e a validação de algoritmos de aprendizado \gls{DRL} voltados para redes veiculares exigem ambientes simulados capazes de replicar a natureza não estacionária e volátil do mundo real. Diferente de aplicações computacionais estáticas, os modelos precisam interagir com um ecossistema que engloba tanto a física da movimentação quanto a complexidade das comunicações sem fio e da carga de processamento~\cite{uddin2025}.

Para modelar a mobilidade com precisão, simuladores de tráfego microscópico, como o \textit{SUMO} (Simulation of Urban MObility)~\cite{lopez2018microscopic} e o \textit{PTV Vissim}~\cite{ptv2020vissim}, são amplamente adotados na literatura~\cite{souza2021, kumar2023}. Ferramentas desse tipo permitem a criação de malhas rodoviárias e urbanas e simulam o comportamento dinâmico dos veículos por meio de modelos matemáticos, capturando variações de velocidade, aceleração e respostas a semáforos~\cite{li2025}. Em paralelo, a dinâmica das comunicações --- que inclui o desvanecimento do canal, perdas de propagação e bloqueios físicos (condições de visada \textit{LoS/NLoS}) --- é frequentemente emulada por simuladores de rede de eventos discretos, destacando-se o ns-3, o qual costuma ser integrado ao SUMO para formar plataformas de co-simulação~\cite{souza2021, uddin2025}.

Em relação ao treinamento do modelo de inteligência artificial, o ambiente de simulação é abstraído matematicamente para que o agente observe os estados e tome decisões. Para viabilizar esse treinamento, trabalhos recentes constroem ambientes customizados baseados em padrões como o \textit{OpenAI Gym} (recentemente sucedido pelo \textit{Gymnasium})~\cite{brockman2016openai}, que gerenciam o ciclo de observação do estado, execução da ação e recebimento da recompensa~\cite{uddin2025} para treinar e otimizar os pesos das redes neurais~\cite{moghaddasi2024, kumar2023, wu2025}.

Por fim, para evitar que o treinamento ocorra de forma puramente teórica, a literatura atual enfatiza o uso de dados do mundo real (\textit{real-world traces}) dentro do ambiente de simulação~\cite{farimani2024}. Algumas abordagens incorporam trajetórias extraídas de frotas reais ou bases de dados de servidores comerciais --- como o \textit{Google Cluster Dataset} --- para emular as flutuações estocásticas na geração de tarefas e a exigência de recursos de \gls{CPU}, Memória e Largura de Banda~\cite{long2025}. Essa imersão de dados reais no ambiente simulado garante que as políticas de descarregamento aprendidas pelo agente \gls{DRL} sejam aplicáveis e estáveis diante das dinâmicas de arquiteturas \gls{VEC} no mundo físico.

% \begin{landscape}
\section*{Discussão}

\begin{table}[h!]
    \centering
    \Caption{\label{tab:estrategia_checkmark}Análise comparativa dos trabalhos relacionados por Função Objetivo, Estado, Recompensa e Ambiente Simulado}
    \setlength{\arrayrulewidth}{0.6pt}
    \renewcommand{\arraystretch}{1.35}
    \resizebox{\textwidth}{!}{
    \begin{tabular}{|l|c|c|c|c|c|c|c|c|c|c|c|c|c|c|c|c|c|}
        \hline
        % ── Linha 1: cabeçalhos de grupo ─────────────────────────────────
        \multirow{3}{*}{Referência}
            & \multicolumn{4}{c|}{Função Objetivo}
            & \multicolumn{4}{c|}{Estado}
            & \multicolumn{5}{c|}{Recompensa}
            & \multicolumn{4}{c|}{Ambiente Simulado} \\ \cline{2-18}
        % ── Linha 2: sub-grupos ───────────────────────────────────────────
        &  &  &  &  &  &  &  &  &
            & \multicolumn{2}{c|}{Pesos}
            & \multicolumn{2}{c|}{Penalização}
            & \multicolumn{2}{c|}{Mobilidade}
            & \multicolumn{2}{c|}{DES} \\ \cline{11-18}
        % ── Linha 3: todos os cabeçalhos folha, rotacionados ─────────────
        & \rotatebox{90}{Latência\hspace{0.3cm}}
            & \rotatebox{90}{Energia\hspace{0.3cm}}
            & \rotatebox{90}{Utilidade\hspace{0.3cm}}
            & \rotatebox{90}{Financeiro\hspace{0.3cm}}
            & \rotatebox{90}{Tarefa\hspace{0.3cm}}
            & \rotatebox{90}{Recursos\hspace{0.3cm}}
            & \rotatebox{90}{Tendências\hspace{0.3cm}}
            & \rotatebox{90}{Escore\hspace{0.3cm}}
            & \rotatebox{90}{Escalonamento\hspace{0.3cm}}
            & \rotatebox{90}{Estáticos\hspace{0.3cm}}
            & \rotatebox{90}{Dinâmicos\hspace{0.3cm}}
            & \rotatebox{90}{Discreta\hspace{0.3cm}}
            & \rotatebox{90}{Contínua\hspace{0.3cm}}
            & \rotatebox{90}{SUMO\hspace{0.3cm}}
            & \rotatebox{90}{Vissim\hspace{0.3cm}}
            & \rotatebox{90}{Customizado\hspace{0.3cm}}
            & \rotatebox{90}{Simulador\hspace{0.3cm}} \\ \hline
        % ── Dados ────────────────────────────────────────────────────────
       ~\cite{baker2024}   & \checkmark & \checkmark & \checkmark &            &            & \checkmark &            & \checkmark & \checkmark & \checkmark &            & \checkmark &            &            &            & \checkmark &            \\ \hline
       ~\cite{wei2024}     & \checkmark & \checkmark &            &            & \checkmark & \checkmark &            &            & \checkmark & \checkmark &            &            & \checkmark &            &            & \checkmark &            \\ \hline
       ~\cite{chen2025}    & \checkmark &            &            &            &            & \checkmark & \checkmark &            &            & \checkmark &            & \checkmark &            &            &            & \checkmark &            \\ \hline
       ~\cite{ma2025}      & \checkmark & \checkmark &            & \checkmark &            & \checkmark &            &            & \checkmark & \checkmark &            & \checkmark &            &            &            & \checkmark &            \\ \hline
       ~\cite{kumar2023}   & \checkmark & \checkmark & \checkmark &            &            & \checkmark & \checkmark &            & \checkmark &            &            &            & \checkmark & \checkmark & \checkmark & \checkmark &            \\ \hline
       ~\cite{kumari2025}  & \checkmark & \checkmark & \checkmark &            & \checkmark & \checkmark &            &            & \checkmark &            & \checkmark & \checkmark &            &            &            & \checkmark &            \\ \hline
       ~\cite{tian2025}    & \checkmark & \checkmark &            &            &            & \checkmark &            &            & \checkmark & \checkmark &            & \checkmark & \checkmark &            &            & \checkmark &            \\ \hline
       ~\cite{bao2025}     & \checkmark & \checkmark &            &            &            & \checkmark &            &            & \checkmark & \checkmark &            &            & \checkmark &            &            & \checkmark &            \\ \hline
       ~\cite{fan2024}     & \checkmark & \checkmark &            &            & \checkmark & \checkmark &            &            &            &            & \checkmark & \checkmark &            &            &            & \checkmark &            \\ \hline
       ~\cite{farimani2024}& \checkmark & \checkmark & \checkmark &            &            &            & \checkmark &            &            &            &            & \checkmark &            &            &            & \checkmark &            \\ \hline
       ~\cite{shi2023}     & \checkmark & \checkmark &            &            &            &            &            &            &            &            & \checkmark &            &            &            &            & \checkmark &            \\ \hline
       ~\cite{vieira2025} & \checkmark & \checkmark &            &            & \checkmark & \checkmark &            & \checkmark &            & \checkmark &            &            &            &            &            & \checkmark &            \\ \hline
    \end{tabular}
    }
    \Fonte{Elaborado pelo autor}
\end{table}
% \end{landscape}

Segundo~\cite{miriyala2026}, apesar dos notáveis avanços proporcionados pela integração do \gls{DRL} no descarregamento de tarefas, a literatura recente aponta limitações que dificultam que os modelos teóricos sejam utilizados no mundo real. Uma análise rigorosa conduzida por~\cite{miriyala2026} revela que a maioria das abordagens atuais sofre de simplificações excessivas na modelagem do ambiente. Frequentemente, os trabalhos utilizam representações de estado baseadas em vetores planos (\textit{flat vectors}), o que ignora as complexas relações estruturais e topológicas entre os veículos e a infraestrutura. Essa lacuna evidencia a necessidade de incorporar arquiteturas diferentes, como as \gls{GNN} para extrair representações topológicas antes da tomada de decisão. Além disso, os ambientes simulados costumam assumir canais de comunicação idealizados e conectividade persistente, negligenciando falhas físicas comuns nas redes veiculares, como quedas intermitentes de conexão (\textit{dropouts})~\cite{miriyala2026}.

Outra lacuna está relacionada a carga computacional imposta pelos próprios modelos de aprendizado de máquina. Grande parte das pesquisas assume que os servidores de borda possuem recursos substanciais de \textit{hardware}, deixando um vácuo no estudo em que dispositivos possuem severas restrições energéticas, de memória e de processamento (teto operacional de \textit{FLOPS}). Para contornar esse obstáculo e viabilizar a implantação,~\cite{miriyala2026} destacam que pesquisas futuras devem focar no desenvolvimento de arquiteturas e representações mais leves, aplicando técnicas de compressão de modelos. 

No que tange à adaptabilidade dos algoritmos, observa-se que as políticas \gls{DRL} são tradicionalmente treinadas e permanecem imutáveis durante a implantação em produção. Contudo, em cenários veiculares não estacionários, em que o ambiente sofre alterações constantes devido a mudanças de distribuição das cargas (\textit{distribution shifts}), tais políticas estáticas divergem e rapidamente se tornam obsoletas e ineficientes~\cite{miriyala2026}. Identifica-se que há espaço na exploração de mecanismos de adaptação \textit{online} de modo a evoluir a política computada pelo modelo.

Por fim, a transição destas propostas para sistemas reais esbarra na falta de avaliação comparativa confiável (\textit{benchmarking}). A disparidade entre os estudos dificulta extrapolações e comparações parciais justas sobre as potenciais eficiências obtidas.~\cite{miriyala2026} criticam o fato de que a literatura tende a se atentar demasiadamente com as constatações de performance baseadas no treinamento ideal dos agentes.

\section{Considerações Finais}
\label{sec:consideracoes-finais-revisao}
Conforme discutido ao longo deste capítulo, as abordagens baseadas em \gls{DRL} têm se solidificado como a direção mais promissora para o gerenciamento de recursos em arquiteturas de borda veicular. A integração com técnicas emergentes demonstra que há flexibilidade suficiente para abarcar os altos desafios impostos pelos cenários veiculares práticos. No entanto, as lacunas destacadas --- como a predominância de simulações com observabilidade perfeita e cenários determinísticos --- salientam a necessidade crítica de propostas que incorporem incerteza, variação de rajadas e carga de trabalho. Este alinhamento direciona o desenvolvimento da solução apresentada nos capítulos seguintes, projetada com vistas a preencher estas falhas e entregar maior robustez.

% \begin{table}[h]
%     \centering \caption{Comparativa de Padrões e Tecnologias de Comunicação} 
%     \begin{tabular}{|l|l|l|l|c|} 
%         \hline \textbf{Artigo} & \textbf{Tecnologia} & \textbf{Servidor} & \textbf{Principal Inovação} & \textbf{Latência} \\ 
%         \hline 
%        ~\cite{Moghaddasi2023} & 5G / \gls{V2I} & BS / Edge & Data Offloading via DDQN & Melhoria de 9,8\% \\ 
%         \hline~\cite{sidwini} & Edge / Cloud & Edge / Cloud & Descarregamento de tarefas Otimizado via \gls{SAC} & Var. \\ 
%         \hline~\cite{li2025} & 6G / \gls{V2X} & MEC (RSU/BS) & Offloading e Associação via \gls{DQN} & Redução de 16,4\% \\ 
%         \hline~\cite{meghamala2025} & 5G / MEC & gNodeB & Alocação de Recursos via MARL & Otimizada \\ 
%         \hline 
%     \end{tabular} 
% \end{table}

% O estudo de~\cite{mondal2025} aborda o desafio de utilizar \gls{DRL}, especificamente um controlador DDQN, para orquestrar de forma inteligente o \textit{handover} e a seleção entre as interfaces \gls{Uu} (celular) e \gls{PC5} (sidelink). 

% Por sua vez,~\cite{heo2026} foca no \textit{streaming} de vídeo interveicular, propondo um framework híbrido que utiliza a via \gls{V2V} como primária e a via \gls{V2N2V} (apoiada por \gls{MEC}) para retransmissões seletivas, garantindo uma latência sob o limite crítico de 50 ms. 


% apresenta o framework \gls{3M}, que integra o Gerenciamento de Mobilidade, tecnologia Multi-RAT e \gls{ML}, buscando garantir um suporte hiper-rápido e ultraconfiável para redes veiculares \gls{6G}, com a projeção de latências que podem chegar a $< 1$ ms para comunicações críticas.
% \begin{table}[h]
% \centering
% \caption{Comparativa das Abordagens de \textit{Descarregamento de Tarefas} Assistidas por MEC}
% \label{tab:comparativa_mec}
% \renewcommand{\arraystretch}{1.3} % Melhora o espaçamento interno das linhas
% \begin{tabular}{|p{2.5cm}|p{3cm}|p{3.5cm}|p{6cm}|}
% \hline
% \textbf{Artigo} & \textbf{Arquitetura} & \textbf{Algoritmos Base} & \textbf{Foco e Uso do \gls{MEC}} \\ \hline

% Zhao et al.~\cite{zhao2025} & \gls{VEC} (\gls{MEC} em \gls{RSU}) & \gls{DQN} + Algoritmo \gls{KM} + \gls{LSTM} & Otimiza o descarregamento de múltiplas tarefas na borda. Utiliza detecção de distância (via predição) para evitar falhas por \textit{handover} frequente entre \gls{RSU}s, focando em minimizar a latência total. \\ \hline

% Wu et al.~\cite{wu2025} & \gls{MEC} (5G/6G) & \gls{GAN} + \gls{FL} & Foca em implantação leve e eficiência energética. Usa o \gls{MEC} para gerenciar \gls{FL} e \gls{GAN} na expansão de amostras, resolvendo o desequilíbrio de dados sem sobrecarregar a rede. \\ \hline

% Mohammad et al.~\cite{mohammad2025} & EMEC (\gls{MEC} para \gls{IoV}) & \gls{DRL} + \gls{PSO} & Apresenta framework EMEC-IoVTOF que delega tarefas para a borda visando reduzir latência e energia. O uso do \gls{PSO} ajuda as decisões de \gls{DRL} a fugirem de armadilhas de ótimos locais. \\ \hline

% Liu et al.~\cite{liu2025} & \gls{C-MEC} (Colaboração Nuvem-Borda) & Aproximação de Bernstein + SCA & Divide o esforço: tarefas sensíveis a atraso vão para \gls{MEC}, e tarefas tolerantes vão para nuvem. Lida com incertezas de canal (mobilidade) e restrições de interferência de forma probabilística. \\ \hline
% \end{tabular}
% \end{table}

% A Tabela~\ref{tab:padrao_comunicacao} sumariza os aspectos de padrão de comunicação e servidor dos trabalhos relacionados mencionados nas subseções anteriores.
% Para contornar limitações de processamento local, tarefas monolíticas estáticas podem ser convertidas utilizando técnicas de divisão de tarefas (\textit{task splitting})~\cite{zhang2024}. Essa conversão viabiliza abordagens de descarregamento parcial, permitindo que cargas de trabalho sejam paralelizadas simultaneamente entre dispositivos locais, servidores \gls{MEC} e infraestrutura de nuvem~\cite{zhang2024, long2025}. No entanto, a implementação prática desse paradigma em sistemas de transporte inteligente enfrenta desafios operacionais severos devido à alta mobilidade dos nós, que provoca flutuações rápidas nos enlaces de comunicação e mudanças extremas na topologia da rede~\cite{kandali2025}.

% Diante de espaços de estado e de ação com alta dimensionalidade e dinâmica não convexa, técnicas tradicionais embasadas em otimização matemática estrita ou teoria dos jogos pura sofrem com baixa escalabilidade e falham em realizar adaptações em tempo real de forma eficiente~\cite{xie2025, zhao2025}. Para superar tais barreiras, soluções contemporâneas precisam transcender esquemas heurísticos simples e focar na otimização conjunta das decisões de descarregamento, escalonamento e alocação de recursos computacionais e de rádio~\cite{zhang2024, qin2020}. Nesse contexto, modelos \gls{DRL} ganham destaque por aprenderem políticas adaptativas baseadas em tentativa e erro e mecanismos de recompensa, adequando-se melhor a ambientes complexos~\cite{zhao2025}.

% Posteriormente, discute-se a adequação dessas abordagens à luz dos desafios expostos.

% O compilado de trabalhos apresentado a seguir categoriza as soluções segundo seus objetivos, engenharia de estado e de recompensa, além dos modelos \gls{DRL} utilizados.

% Apesar do sucesso alcançado por arquiteturas \gls{DRL} na orquestração de recursos, uma lacuna crítica prevalece em grande parte da literatura atual: a omissão das incertezas inerentes à mobilidade e ao estado do sistema, o que muitas vezes reduz o problema a um cenário de observabilidade perfeita~\cite{liu2025}. Muitos modelos assumem erroneamente conhecimento exato e determinístico das condições do canal e da disponibilidade de recursos no exato momento da decisão~\cite{liu2025}. Na prática, a imprevisibilidade do processo de descarregamento na borda~\cite{abbas2025}, aliada à alta volatilidade e ao atraso natural de telemetria, faz com que decisões baseadas em premissas estáticas sejam ineficientes~\cite{kandali2025, zhao2025}. Tais fatores podem levar agentes a tomar decisões desatualizadas, resultando em contenção de recursos, sobrecarga em servidores específicos e falhas no atendimento dos requisitos sob rajadas de tráfego~\cite{kandali2025, abbas2025}.

% Para garantir decisões mais robustas e estabilidade de convergência, a pesquisa atual enfatiza a necessidade de adotar mecanismos avançados de abstração. A introdução de técnicas que modelem a incerteza do ambiente — como restrições probabilísticas, inferência baseada em lógica difusa (Fuzzy) ou algoritmos metaheurísticos híbridos — pode enriquecer os modelos \gls{DRL}. Tais integrações auxiliam na quantificação de variáveis incertas e aceleram a descoberta de soluções globais e confiáveis para o descarregamento de tarefas, superando as restrições da observabilidade puramente determinística~\cite{liu2025, zhao2025}.



% Tais algoritmos permitem que o sistema aprenda políticas ótimas de orquestração por meio da interação contínua com o ambiente. Trabalhos recentes expandem essa fronteira combinando abordagens \gls{DRL} com Aprendizado por Transferência (\textit{Transfer Learning}), acelerando drasticamente a adaptação de nós recém-conectados e reduzindo o tempo de treinamento~\cite{baker2024}. Outra estratégia inovadora é a integração com Gêmeos Digitais (Cybertwin) em redes \gls{6G}, permitindo que o treinamento e as simulações de orquestração ocorram virtualmente antes da execução real no ambiente veicular~\cite{wei2024}. Além disso, arquiteturas distribuídas baseadas em Aprendizado Federado (FRL) ganham força como estratégia para permitir a colaboração entre múltiplos veículos na formulação de políticas sem a necessidade de centralizar dados brutos ou expor informações privadas~\cite{chen2025}.

% No que tange à função de recompensa, existem abordagens que adotam pesos dinâmicos adaptativos que se ajustam em tempo real de acordo com as condições momentâneas da rede e com o nível de urgência das tarefas~\cite{kumari2025, fan2024}. A penalização estrita por falhas na execução --- como a violação do tempo máximo tolerável (deadline) --- é ativamente integrada à recompensa para guiar o agente a priorizar o cumprimento dos requisitos de latência~\cite{farimani2024}. Em cenários de descarregamento parcial, a recompensa é formulada dividindo a ação em múltiplos estágios sequenciais, avaliando o impacto do particionamento da tarefa de forma independente da alocação de recursos~\cite{tian2025}. Além disso, a teoria de Lyapunov é frequentemente injetada na recompensa por meio de funções \textit{drift-plus-penalty}, garantindo matematicamente que o agente seja penalizado caso a redução excessiva de consumo energético coloque em risco a estabilidade das filas do sistema a longo prazo~\cite{kumar2023}.


% \begin{table}[ht!]
%     \centering
%     \caption{Classificação dos trabalhos relacionados baseada na taxonomia de Estratégia e Ambiente}
%     \label{tab:taxonomia_estrategia_ambiente}
%     \renewcommand{\arraystretch}{1.3}
%     \resizebox{\textwidth}{!}{
%     \begin{tabular}{|l|p{3.8cm}|p{2.5cm}|p{2.5cm}|p{2.8cm}|p{3.5cm}|}
%         \hline
%         \multirow{2}{*}{\textbf{Referência}} & \multicolumn{4}{c|}{\textbf{Estratégia}} & \multirow{2}{*}{\textbf{Ambiente}} \\ \cline{2-5}
%         & \textbf{Objetivo} & \textbf{Eng. de Estado} & \textbf{Eng. de Recompensa} & \textbf{Modelos} & \\ \hline
        
%        ~\cite{baker2024} & $-$ Tempo de Resposta \newline $-$ Cons. Energético \newline $+$ Utilidade do Sistema & Latência, Buffer, CPU & Latência, Energia & \gls{DQN} + \textit{Transfer Learning} & Estocástico, Carga de Trabalho \\ \hline
        
%        ~\cite{wei2024} & $-$ Tempo de Resposta \newline $-$ Cons. Energético \newline $-$ Sobrecarga & Canal, Latência, CPU & Função Custo & QDPG (DDPG + D3QN) & Estocástico, Carga de Trabalho \\ \hline
        
%        ~\cite{chen2025} & $-$ Tempo de Resposta & Canal, Buffer, CPU & Latência & \gls{DQN} Federado & Estocástico, Carga de Trabalho \\ \hline
        
%        ~\cite{ma2025} & $-$ Tempo de Resposta \newline $-$ Cons. Energético & Canal, Buffer, CPU & Função Custo & P-DQN & Estocástico \\ \hline
        
%        ~\cite{kumar2023} & $-$ Tempo de Resposta \newline $-$ Cons. Energético \newline $-$ Sobrecarga & Canal, Latência, Buffer, CPU & Função Custo & L-MADDPG (Ator-Crítico) & Estocástico, Incerteza, Carga de Trabalho \\ \hline
        
%        ~\cite{kumari2025} & $-$ Tempo de Resposta \newline $-$ Cons. Energético \newline $+$ Utilidade do Sistema & Canal, Latência, Buffer, CPU & Função Custo & \gls{DQN} & Estocástico \\ \hline
        
%        ~\cite{tian2025} & $-$ Tempo de Resposta \newline $-$ Cons. Energético & Canal, Buffer, CPU & Latência, Energia & SIDDPG (Ator-Crítico) & Estocástico, Carga de Trabalho \\ \hline
        
%        ~\cite{bao2025} & $-$ Tempo de Resposta \newline $-$ Cons. Energético & Canal, Latência, Buffer, CPU & Função Custo & \gls{SAC} & Estocástico, Incerteza \\ \hline
        
%     \end{tabular}
%     }
% \end{table}